\documentclass[journal]{IEEEtran}

% 中文支持（必须在最前面）
\usepackage{xeCJK}
\setCJKmainfont{SimSun}  % 或使用 "Microsoft YaHei", "FandolSong" 等
\setCJKsansfont{SimHei}

% 基础数学包
\usepackage{amsmath,amsfonts}

% 算法伪代码
\usepackage{algorithmic}
\usepackage{algorithm}

% 图形和表格
\usepackage{graphicx}
\usepackage{booktabs}     % 专业表格线条（\toprule, \midrule, \bottomrule）
\usepackage{multirow}     % 表格多行合并
\usepackage{array}        % 增强表格功能
\usepackage{adjustbox}    % 调整表格宽度

% 参考文献
\usepackage{cite}

% 断词规则
\hyphenation{op-tical net-works semi-conduc-tor IEEE-Xplore}

% 图片路径
\graphicspath{{figures/}}

% 超链接
\usepackage[unicode]{hyperref}

\begin{document}

% ========== 标题和作者 ==========
\title{技术路线调研报告}

% 作者信息留空，由用户自行填写
\author{}

\maketitle

% ========== 摘要 ==========
\begin{abstract}
本报告针对\textbf{[研究问题]}进行了系统性的技术路线调研。通过分析XX篇相关文献，梳理出N类主要技术路线：[简要列出各类方法]。报告从理论基础、核心优势、局限性、适用范围等多个维度对各类技术路线进行了深入分析和对比，并提出了[问题]的改进方向。
\end{abstract}

% ========== 关键词 ==========
\begin{IEEEkeywords}
技术路线, [关键词1], [关键词2], [关键词3], [关键词4]
\end{IEEEkeywords}

% ========== 第一部分：引言 ==========
\section{引言}

\subsection{问题背景}
[描述研究问题的背景、重要性、应用场景]

\subsection{技术背景}
[介绍相关领域的总体技术状况]

\subsection{报告结构}
[简要说明报告各章节内容]

% ========== 第二部分：技术路线分析 ==========
\section{技术路线分析}

\subsection{概述}
[简要说明识别出的技术路线类别及其数量]

% ----- 技术路线类别1 -----
\subsection{[类别1名称]方法}

\subsubsection{方法原理}
[描述该类方法的核心思想，必须引用参考文献]

该类方法基于[理论基础]，其核心思想是[简述]。关键技术包括：
\begin{itemize}
    \item [技术点1] \cite{ref1, ref2}
    \item [技术点2] \cite{ref3}
    \item [技术点3] \cite{ref4}
\end{itemize}

[如适用，添加数学模型或算法描述]

\subsubsection{代表性工作}
[分析2-3篇代表性文献]

\textbf{[作者等]} 提出了[方法名称] \cite{ref1}，其主要贡献是：
\begin{itemize}
    \item [贡献1]
    \item [贡献2]
    \item [贡献3]
\end{itemize}

\textbf{[作者等]} 在上述工作基础上进行了改进 \cite{ref2}，引入了[改进点]。

\subsubsection{演进脉络}
[按时间梳理该技术路线的发展过程，指出关键突破]

该技术路线的发展可分为以下阶段：
\begin{enumerate}
    \item [阶段1，时间范围]：[特点] \cite{ref1}
    \item [阶段2，时间范围]：[特点] \cite{ref2, ref3}
    \item [阶段3，当前]：[特点] \cite{ref4}
\end{enumerate}

\subsubsection{批判性分析}
\textbf{假设条件}：[说明方法的假设前提及其合理性]

\textbf{核心优势}：
\begin{itemize}
    \item [优势1]
    \item [优势2]
\end{itemize}

\textbf{主要局限}：
\begin{itemize}
    \item [局限1]
    \item [局限2]
\end{itemize}

\textbf{适用范围}：[说明适用的场景和边界条件]

\textbf{实际挑战}：[说明工程实现和部署中的难点]

% ----- 技术路线类别2 -----
\subsection{[类别2名称]方法}
[使用与上述相同的结构]

% ----- 技术路线类别3 -----
\subsection{[类别3名称]方法}
[使用与上述相同的结构]

% ========== 第三部分：技术路线对比 ==========
\section{技术路线对比}

\subsection{特征对比}
使用以下表格对比各类技术路线的特征（所有描述必须基于文献内容并引用）：

\begin{table*}[htbp]
\centering
\caption{各技术路线特征对比}
\begin{tabular}{lcccc}
\toprule
技术路线 & 理论基础 & 核心优势 & 主要局限 & 适用场景 \\
\midrule
[类别1]方法 \cite{ref1,ref2} & [理论基础] & [优势] & [局限] & [场景] \\
[类别2]方法 \cite{ref3} & [理论基础] & [优势] & [局限] & [场景] \\
[类别3]方法 \cite{ref4,ref5} & [理论基础] & [优势] & [局限] & [场景] \\
\bottomrule
\end{tabular}
\end{table*}

\subsection{性能指标对比}
使用以下表格对比各类技术路线的性能指标（所有数据必须来自文献并注明出处）：

\begin{table*}[htbp]
\centering
\caption{性能指标对比}
\begin{tabular}{lcccc}
\toprule
技术路线 & [指标1名称] & [指标2名称] & [指标3名称] & 计算复杂度 \\
\midrule
[类别1]方法 \cite{ref1} & [数据] & [数据] & [数据] & [复杂度] \\
[类别2]方法 \cite{ref3} & [数据] & [数据] & [数据] & [复杂度] \\
[类别3]方法 \cite{ref4} & [数据] & [数据] & [数据] & [复杂度] \\
\bottomrule
\end{tabular}
\label{tab:performance}
\end{table*}

\subsection{综合评估}
[基于加权排序的结果，总结各类技术路线的适用性]

根据[研究问题]的特点，我们确定了以下指标权重：[列出权重]。综合评分结果为：[排名结果]。

% ========== 第四部分：总结与展望 ==========
\section{总结与展望}

\subsection{问题解决现状}
[总结研究问题的解决状况]

通过对各类技术路线的分析，我们认为[研究问题]：
\begin{itemize}
    \item [已解决的部分]：[说明哪些方面已有成熟方案]
    \item [部分解决的部分]：[说明哪些方面有方案但不完善]
    \item [未解决的部分]：[说明哪些方面仍是开放问题]
\end{itemize}

\subsection{改进方向}
[基于前述分析的局限性，提出具体的改进建议]

\subsubsection{技术融合方向}
[说明不同技术路线融合的可能性]

\subsubsection{新兴技术应用}
[说明有潜力的新兴技术如何应用于该问题]

\subsubsection{具体改进建议}
[提出2-3条具体的、可操作的改进方向，每条都应该基于前述分析的某类技术的局限性]

\begin{enumerate}
    \item [改进方向1]：针对[某类技术]的[某局限性]，可以[具体建议]。
    \item [改进方向2]：结合[技术A]和[技术B]的优势，有望[预期效果]。
    \item [改进方向3]：[基于某技术路线的不足]提出的改进方向。
\end{enumerate}

% ========== 参考文献 ==========
\bibliographystyle{IEEEtran}
\bibliography{myref}

\end{document}
